% Scénarios de test consolidés — chapitre 5 (workflow transverse)
\small
\PFETableSetup
\setlength{\LTpre}{0pt}
\setlength{\LTpost}{0pt}
\setlength{\LTleft}{\fill}
\setlength{\LTright}{\fill}
\setlength{\tabcolsep}{3.5pt}
\renewcommand{\arraystretch}{1.25}

\begin{longtable}{|>{\centering\arraybackslash}m{2.8cm}|
>{\raggedright\arraybackslash}p{2.5cm}|
>{\raggedright\arraybackslash}p{3.8cm}|
>{\raggedright\arraybackslash}p{4.3cm}|
>{\centering\arraybackslash}p{1.1cm}|}
\caption{Synthèse des cas de test du sprint 3}\label{tab:ch5-scenarios-tests}\\
\hline
\tableHeaderStyle
\multicolumn{1}{|c|}{\tableHeaderCell{Nom du cas d'utilisation}} &
\multicolumn{1}{c|}{\tableHeaderCell{Cas de test}} &
\multicolumn{1}{c|}{\tableHeaderCell{Description du cas}} &
\multicolumn{1}{c|}{\tableHeaderCell{Résultat attendu}} &
\multicolumn{1}{c|}{\tableHeaderCell{Résultat réel}} \\
\hline
\endfirsthead
\hline
\multicolumn{5}{|c|}{\tablename\ \thetable{} --- \textit{(suite)}} \\
\hline
\tableHeaderStyle
\multicolumn{1}{|c|}{\tableHeaderCell{Nom du cas d'utilisation}} &
\multicolumn{1}{c|}{\tableHeaderCell{Cas de test}} &
\multicolumn{1}{c|}{\tableHeaderCell{Description du cas}} &
\multicolumn{1}{c|}{\tableHeaderCell{Résultat attendu}} &
\multicolumn{1}{c|}{\tableHeaderCell{Résultat réel}} \\
\hline
\endhead
\hline
\multicolumn{5}{|r|}{\textit{Suite page suivante\ldots}} \\
\hline
\endfoot
\hline
\endlastfoot

\needspace{18\baselineskip}
\multirow{3}{=}{\centering Valider un plan CSR} &
Validation corporate (mode direct) &
Un utilisateur Corporate valide un plan soumis lorsque le mode le permet sans étape site préalable. &
Le plan passe au statut validé et n'est plus modifiable. &
Vrai \\
\cline{2-5}
 &
Validation en deux étapes &
Le Corporate tente de valider avant la validation du site lorsque le mode l'exige. &
Le système refuse et indique que la validation du site est requise en premier. &
Vrai \\
\cline{2-5}
 &
Consultation liste des plans soumis &
Ouverture de l'interface listant les plans en attente de décision. &
Les informations essentielles (site, année, statut, budget) sont affichées. &
Vrai \\
\hline

\needspace{12\baselineskip}
\multirow{2}{=}{\centering Rejeter un plan CSR} &
Motif obligatoire manquant &
Tentative de rejet sans saisir le motif de rejet. &
La soumission est refusée et le champ motif est signalé comme obligatoire. &
Vrai \\
\cline{2-5}
 &
Rejet avec motif valide &
Saisie d'un motif de rejet et confirmation du rejet. &
Le plan est rejeté, le statut est mis à jour et le motif est conservé. &
Vrai \\
\hline

\needspace{18\baselineskip}
\multirow{3}{=}{\centering Soumettre une demande\\de modification pour un plan} &
Données incomplètes &
Soumission d'une demande sans justification ou sans durée d'ouverture requise. &
Message d'erreur listant les champs à compléter. &
Vrai \\
\cline{2-5}
 &
Pièce jointe non conforme &
Ajout d'un fichier de type ou taille non autorisé. &
Le fichier est refusé avec un message explicite. &
Vrai \\
\cline{2-5}
 &
Soumission valide &
Demande complète sur un plan verrouillé avec justification et durée. &
La demande est enregistrée avec le statut « en attente » et associée au plan. &
Vrai \\
\hline

\needspace{12\baselineskip}
\multirow{2}{=}{\centering Traiter une demande\\de modification} &
Approbation corporate &
Examen puis approbation d'une demande en attente. &
Le statut de la demande est mis à jour et le site est notifié. &
Vrai \\
\cline{2-5}
 &
Rejet avec justification &
Rejet d'une demande avec commentaire obligatoire. &
Le statut est mis à jour et le motif est conservé pour le site. &
Vrai \\
\hline

\needspace{12\baselineskip}
\multirow{2}{=}{\centering Recevoir des\\notifications} &
Notification après validation &
Un événement de validation génère une notification pour les acteurs concernés. &
La notification apparaît dans le centre de notifications. &
Vrai \\
\cline{2-5}
 &
Marquer comme lue &
L'utilisateur ouvre ou marque une notification comme lue. &
L'état de lecture est mis à jour et le compteur reflète le changement. &
Vrai \\
\hline

\needspace{12\baselineskip}
\multirow{2}{=}{\centering Gérer les tâches} &
Tâche créée après un événement &
Soumission ou validation déclenchant la création d'une tâche assignée. &
La tâche apparaît dans la liste des tâches de l'utilisateur concerné. &
Vrai \\
\cline{2-5}
 &
Ouverture d'une tâche &
Clic sur une tâche pour accéder au contexte (plan ou demande). &
La navigation mène à l'écran attendu. &
Vrai \\
\hline

\needspace{12\baselineskip}
\multirow{2}{=}{\centering Consulter le\\journal d'audit} &
Affichage chronologique &
Consultation de la liste des événements par un profil Corporate autorisé. &
Les actions sont listées avec date, type, entité et auteur. &
Vrai \\
\cline{2-5}
 &
Filtrage par critères &
Application de filtres (site, type d'action, période). &
La liste se restreint aux enregistrements correspondants. &
Vrai \\
\hline

\end{longtable}
\normalsize
